% Printed212; principalPDF219.
\par\clearpage
\noindent212\hfill\textsc{Briefwechsel.}\hfill\mbox{}\par\medskip
\begin{center}Concise Musterrechnung.\\Neuwerk.\end{center}
\[
\begin{array}{r@{\qquad}r@{\qquad}r}
-266575{,}038&+95076{,}254&\\[8pt]
\begin{array}[t]{r}
133{,}287\,519\\10{,}663\,002\\\hline
143{,}950\,521\\2^\circ\,23'\,57''{,}03\\51\quad23\quad20{,}61\\\hline
53\quad47\quad17{,}64\\8\quad15{,}73\\\hline
53\quad55\quad33{,}37\\-31{,}40\\\hline
53^\circ\,55'\quad1''{,}97\\\hline
\end{array}&
\begin{array}[t]{r}
4{,}978\,0721\\8{,}508\,8620\\\hline
3{,}486\,9341\\9{,}769\,9904\\0{,}137\,5588\\9{,}956\,1442\\1{,}403\,2842\\\hline
3{,}716\,9437\\-764\\\hline
1{,}496\,9872\\-850\\\hline
3{,}624\,4929\\-923
\end{array}&
\begin{array}[t]{r}
7{,}43389\\\hline5{,}44934\\5{,}49566\\5{,}53146\\[8pt]
5210''{,}355\\1^\circ\,26'\,50''{,}355\\\hline
4211''{,}149\\1^\circ\,10'\,11''{,}15\\\hline
\end{array}
\end{array}
\]
\bigskip
\begin{center}\textsc{Gauss} an \textsc{Schumacher}.\quad Göttingen, 30. April 1830.\end{center}
. . . . . . Ich fahre jetzt fort mit dem, was die Hülfstafeln betrifft; heute nur die Formeln für die übrigen drei Columnen, die bei Berechnung der Länge, Breite und Convergenz gebraucht werden, ich weiss aber nicht, ob ich sie in meinem vorigen Briefe mit \(D,E\) und \(F\) bezeichnet habe.

Zahlen, deren Logarithmen angesetzt werden:
\[
\begin{array}{@{}l@{\;}l@{}}
\text{Länge . . . . }[D=]&h\left(\frac34-\frac14\cos2\varphi+\frac{ee}{2(1-ee)}\cos\varphi^4\right)\\[6pt]
\text{Breite . . . . }[E=]&h\left(\frac34+\frac{2-11ee}{4(1-ee)}\cos\varphi^2+\frac{5ee}{2(1-ee)}\cos\varphi^4-\frac{e^4}{(1-ee)^2}\cos\varphi^6\right)\\[6pt]
\text{Convergenz }[F=]&h\left(1-\frac{ee}{1-ee}\cos\varphi^4-\frac{2e^4}{(1-ee)^2}\cos\varphi^6\right).
\end{array}
\]
\[
\text{Hier ist }h=\frac{k\,.\,10^7}{3\,.\,(206265)^2};\qquad
\begin{array}{l}
k=\text{Modulus}=0{,}434\rlap{\raisebox{-.5ex}{\tiny.}}^{\prime}2945\\
\log h=5{,}531\,8128-10.
\end{array}
\]

% Printed213; principalPDF220.
\par\clearpage
\noindent\mbox{}\hfill{\small\textsc{Über die Formeln für die Hannoversche Landesvermessung.}}\hfill213\par\medskip
Es ist zu bemerken, dass diese Formeln vollständig sind, d. i. es sind keine unendliche Reihen, sondern nur diese Glieder. . . . . . .

Um Ihr Vertrauen zu \textsc{Schmidts} Rechnung zu vergrössern, bemerke ich, dass er die zwei Hauptelemente der Erddimensionen viermal berechnet hat, aber nur Einmal hat er wegen Rechnungsfehlers von neuem gerechnet. Nemlich
\begin{list}{}{\setlength{\leftmargin}{2em}\setlength{\itemsep}{0pt}\setlength{\parsep}{0pt}}
\item[1)] Zahlen in meiner Breitenbestimmung etc.

Diese hatten einen Rechnungsfehler enthalten, den er später verbesserte; daher
\item[2)] die Zahlen in seiner Geographie und in Ihren A. N.

Erst später machte ich ihn aufmerksam auf die Correction der in Ostindien gebrauchten Maassstäbe; daher die
\item[3)] Rechnung, deren Resultat in der Vorrede des Buchs.

Endlich hat er seitdem eine vierte Rechnung gemacht, nicht wegen eines Rechnungsfehlers, sondern um die ihm erst nachher bekannt gewordenen Resultate von \textsc{Struves} Gradmessung mit unter die Data aufzunehmen. Das Resultat
\item[4)] ist mir von ihm handschriftlich mitgetheilt und dasselbe, was meinen neuen Hülfstafeln zum Grunde liegt, nemlich
\end{list}
\[
\text{Abplattung}=\frac1{297{,}732};\qquad
\frac{\text{Erdquadrant}}{90}=57008{,}551.
\]
. . . . . .
\bigskip
\begin{center}\textsc{Gauss} an \textsc{Schumacher}.\quad Göttingen, 17. Mai 1831.\end{center}
. . . . . . Es gibt \textls[120]{mehrere} Wege, um aus den Breiten und Längen die Coordinaten zu berechnen, die jeder an seiner Stelle ihre eigenthümlichen Vorzüge haben; für den Fall, wozu Sie solcher Rechnungen bedürfen, ist es am bequemsten, sich der Reihen zu bedienen, die dann so schnell convergiren, dass sehr wenige Glieder hinreichen.

Entsprechen die Coordinaten \(x,y\) der Breite \(\varphi\) und Länge \(\lambda\), so ist die Form diese\textsuperscript{[*)]}:
\[
\begin{array}{l}
x=A-A'\lambda\lambda-A''\lambda^4-\text{etc.}\\
y=B\lambda+B'\lambda^3+\text{etc.},
\end{array}
\]
\par\noindent\rule{4cm}{.4pt}\par
{\footnotesize\noindent\hspace*{\parindent}[*) Vergl. Art. 13, S. 158.]\par}

% Printed214; principalPDF221.
\par\clearpage
\noindent214\hfill\textsc{Briefwechsel.}\hfill\mbox{}\par\medskip
\noindent wo \(A,A',A''\), etc.; \(B,B'\), etc. Functionen von \(\varphi\) sind, deren numerische Werthe man für diejenigen runden Grade (oder halben Grade etc.) zu berechnen hat, welche innerhalb der Karte vorkommen.

Zur Berechnung von \(A\) bedürfen Sie meiner Anleitung nicht, da sie lediglich von der Rectification eines elliptischen Bogens abhängt. Ist nemlich \(\varphi^\circ\) die Polhöhe desjenigen Orts, von wo an man die Coordinaten \(x\) südlich zählt; \(u^\circ\) dessen wirkliche Distanz vom Äquator und \(u\) indefinite die Äquatordistanz des Parallelkreises \(\varphi\) (beide, \(u,u^\circ\), in derjenigen Einheit ausgedrückt, die man für die Coordinaten gewählt hat), so ist \(A=u^\circ-u\). Sie können dazu auch eine Ihnen schon früher mitgetheilte Hülfstafel benutzen\textsuperscript{[*)]}. Die Werthe der andern Coefficienten können Sie nach folgenden Formeln berechnen, die \textls[120]{absolut} genau sind.
\[
\begin{aligned}
A'&=\frac{asc}{2p}\\[4pt]
A''&=\frac{asc}{24(1-ee)^2p}\bigl(5-ee-(6+6ee)ss+(9ee+3e^4)s^4-4e^4s^6\bigr)\\[4pt]
B&=\frac{ac}{p}\\[4pt]
B'&=\frac{ac}{6(1-ee)p}(1-2ss+ees^4).
\end{aligned}
\]
Wobei folgendes zu bemerken ist.

Es bedeuten:
\begin{quote}
\(a\) den Halbmesser des Äquators, \(e\) die Excentricität;
\[
\left.\begin{array}{l}
c=\cos\varphi\\s=\sin\varphi\\p=\surd(1-eess)
\end{array}\right\}\ \text{Kürze halber;}
\]
\end{quote}
\noindent ferner ist für \(\lambda\) der Bogen \(57^\circ17'45''\) als Einheit angenommen: will man also etwa zuletzt die Coordinaten von Grad zu Grad für \(\lambda\) berechnen, so wird man wohl thun, gleich anfangs dem Coefficienten \(B\) den Factor \(\frac1{57{,}296}=\frac{3600}{206265}=\frac\pi{180}\) beizufügen und ebenso die Coefficienten \(A',B',A''\) sogleich mit der zweiten, dritten, vierten Potenz von \(\frac\pi{180}\) zu multipliciren.

Die Formeln fallen etwas einfacher aus, wenn man \(\frac{ee}{1-ee}=m\) setzt und jene danach umschmelzt. In dieser Form habe ich sie selbst zu meinen
\par\noindent\rule{4cm}{.4pt}\par
{\footnotesize\noindent\hspace*{\parindent}[*) Seite 206.]\par}

% Printed215; principalPDF222.
\par\clearpage
\noindent\mbox{}\hfill{\small\textsc{Über die Formeln für die Hannoversche Landesvermessung.}}\hfill215\par\medskip
\noindent Rechnungen angewandt, allein nicht aufgehoben\textsuperscript{[*)]}; ich überlasse also die sehr kleine Arbeit jener Umformung Ihnen selbst. Indem ich \(\varphi^\circ=51^\circ31'47''{,}85\) annehme, für die Abplattung Herrn \textsc{Schmidts} letzte Bestimmung zum Grunde lege und zur Lineareinheit den \(10\,000\,000.\) Theil des Erdquadranten in dieser Gestalt wähle, finde ich, \(\lambda=n\) Grad gesetzt:
\[
\begin{array}{c@{\quad}l@{\qquad}l}
\varphi&x=&y=\\
51^\circ&+58947{,}1-475{,}95nn-0{,}0167n^4&70180{,}0n-0{,}737n^3\\
52^\circ&-52287{,}0-472{,}16nn-0{,}0154n^4&68660{,}7n-0{,}840n^3\\
53^\circ&-163539{,}8-467{,}79nn-0{,}0140n^4&67120{,}2n-0{,}936n^3\\
54^\circ&-274811{,}2-462{,}85nn-0{,}0127n^4&65559{,}1n-1{,}026n^3\\
55^\circ&-386100{,}9-457{,}34nn-0{,}0114n^4&63977{,}8n-1{,}109n^3.
\end{array}
\]
. . . . . .
\bigskip
\begin{center}\textsc{Gauss} an \textsc{Schumacher}.\quad Göttingen, 25. Junius 1831.\end{center}
. . . . . . Was die Berechnung der Coordinaten betrifft, so kommt es auf zwei Aufgaben an, nemlich:

1) Aus der wirklichen Länge \(R\) einer kürzesten Linie auf dem Sphäroid, deren Endpunkte in der Darstellung resp. die Coordinaten \(x,y;\ x',y'\) haben, die Entfernung [in] der Darstellung, d. i. die Grösse \(r=\surd((x'-x)^2+(y'-y)^2)\) zu finden. Hier ist die Auflösung der umgekehrten Aufgabe. Es ist für alle Ihre Fälle mit hinreichender Genauigkeit
\[
\log R=\log r-k\,.\,\frac{(1-ee\sin\varphi^2)^2}{2aa(1-ee)}\,.\,\frac{yy+yy'+y'y'}3.
\]
Da man dabei \(y\) und \(y'\) nicht sehr genau zu kennen braucht, so ist dazu die vorläufige Berechnung der Coordinaten, die man so macht, als ob alles in plano wäre, zureichend, und so dient die Formel auch für die Aufgabe, \(r\) aus \(R\) zu finden. \(k\) ist der Modulus der briggischen Logarithmen.

Zur wirklichen Berechnung setzte ich die Formel in folgende Gestalt:
\[
\log r=\log R+\{\alpha(y+y')^2+\beta(y-y')^2\}q,
\]
\noindent wo
\[
\alpha=\frac{k}{4\,.\,206265},\quad\beta=\frac{k}{12\,.\,206265},\quad q=206265\,.\,\frac{(1-ee\sin\varphi^2)^2}{2aa(1-ee)};
\]
\par\noindent\rule{4cm}{.4pt}\par
{\footnotesize\noindent\hspace*{\parindent}[*) Diese Formeln findet man im Art. 13, S. 157 unten.]\par}

% Printed216; principalPDF223.
\par\clearpage
\noindent216\hfill\textsc{Briefwechsel.}\hfill\mbox{}\par\medskip
\noindent hier ist
\[
\begin{array}{l}
\log\alpha=3{,}72130-10\\
\log\beta=3{,}24418-10.
\end{array}
\]
Für \(q\) habe ich eine Hülfstafel, die bloss von \(x\) abzuhängen braucht, und die, wie ich glaube, ich Ihnen schon einmal mitgetheilt habe\textsuperscript{[*)]}. Offenbar sind \(\alpha,\beta\) von der Maasseinheit unabhängig, aber nicht \(q\), wofür eine Hülfstafel nicht bloss von der Maasseinheit, sondern auch von dem Anfang der \(x\) abhängig ist. Die meinige müssten Sie also, wenn Sie eine andere Einheit und einen andern Anfangspunkt brauchen, erst transformiren.

Übrigens erhalten Sie so \(\log r-\log R\) als Decimalbruch; wollen Sie den Werth gleich in Einheiten der 7. Decimale haben, so brauchen Sie nur \(\log\alpha=0{,}72130,\ \log\beta=0{,}24418\) anzuwenden.

2) Aus einem im Punkte \(P\) gemessenen Winkel zwischen \(PP'\) und \(PP''\) den Winkel zu finden, welcher ihm in der Darstellung correspondirt. Auch hier tritt wieder die umgekehrte Aufgabe an die Stelle.
\begin{center}
\begin{tikzpicture}[x=1cm,y=1cm,font=\normalsize]
\node at (1.8,2) {Sphäroid};\node at (9,2) {Zeichnung in plano};
\draw (0,.7) circle(.035);\node[above left] at (0,.7) {$P''$};
\draw (2.1,1.2) circle(.035);\node[above right] at (2.1,1.2) {$P$};
\draw (2.85,.25) circle(.035);\node[above right] at (2.85,.25) {$P'$};
\draw (7.1,.7) circle(.035);\node[above left] at (7.1,.7) {$p''$};
\draw (9.3,1.2) circle(.035);\node[above right] at (9.3,1.2) {$p$};
\draw (10.05,.25) circle(.035);\node[above right] at (10.05,.25) {$p'$};
\end{tikzpicture}
\end{center}
Es seien die Coordinaten von \(p,p',p''\) respective
\[
\begin{array}{ccc}x,&x',&x'',\\y,&y',&y'',\end{array}
\]
\noindent dann ist:
\[
\begin{aligned}
P'PP''={}&p'pp''-q^*(x'-x)\left(\frac{2y+y'}3\right)\\
&+q^{**}(x''-x)\left(\frac{2y+y''}3\right).
\end{aligned}
\quad\text{(gibt das Gesuchte in Secunden)}
\]
Sie können für \(q^*,q^{**}\) den Werth von \(q\) anwenden, welcher dem Argument \(x\) correspondirt; wollen Sie genauer gehen, so ist für \(q^*\) das Argument \(\frac{2x+x'}3=x+\frac13(x'-x)\) und für \(q^{**}\) das Argument \(\frac{2x+x''}3=x+\frac13(x''-x)\). Ich selbst setze immer \(\frac{2y+y'}3\) in die Form \(y+\frac13(y'-y)\), etc. Sie sehen, dass dann das ganze Verfahren darauf hinausläuft, erst jedem gemessenen Azimuth (von
\par\noindent\rule{4cm}{.4pt}\par
{\footnotesize\noindent\hspace*{\parindent}[*) Seite 210, wo \(H=q\) ist.]\par}

% Printed217; principalPDF224.
\par\clearpage
\noindent\mbox{}\hfill{\small\textsc{Über die Formeln für die Hannoversche Landesvermessung.}}\hfill217\par\medskip
\noindent \(P\) nach \(P'\)) die Correction \(-q(x'-x)(y+\frac13(y'-y))\) beizufügen, woraus das entsteht, was ich Azimuth in plano nenne. Unter Azimuth auf dem Sphäroid verstehe ich hier aber nicht das astronomische Azimuth, d. i. nicht den Winkel mit dem wirklichen Meridian, sondern mit einer Linie, die dem Fundamentalmeridian parallel ist, oder strenger, mit einer Linie auf dem Sphäroid, die in der Darstellung in plano eine Parallele mit der Abscissenlinie gibt.

Azimuth in plano ist also immer \(\operatorname{arc\,tang}\frac{y'-y}{x'-x}\).

Nachdem jene Correctionen angewandt sind, so hat man mit allen Winkeln so zu rechnen, als ob alles in plano wäre, und wie dann die Coordinaten zu berechnen sind, darüber bedürfen Sie keiner Vorschriften. Offenbar ist auch für diese Rechnung eine genäherte Kenntniss von \(x,y,x'\), etc. hinreichend, wie man sie erhält, wenn man anfangs ohne Correction rechnet; will man alles in den \(0''{,}001\) harmonisch haben, so kann man allenfalls die Rechnung, nachdem \(x,y\), etc. schärfer bekannt sind, retouchiren.

Übrigens sehen Sie leicht, dass die obige Rechnung 1), d. i. Übergang von \(R\) auf \(r\), nur bei Einer Linie zu machen ist (der Basis), nachher hat die Kenntniss der einzelnen Linien auf dem Sphäroid, allgemein zu reden, kein Interesse, es sei denn, dass man wieder zu einer neuen Basis gelangt, wo da der verkehrte Weg (von \(r\) nach \(R\)) anzuwenden ist. Ohnehin ist in allen Ihren Fällen \(R\) und \(r\) immer sehr nahe gleich, jedenfalls hat der Unterschied auf das Centriren der Winkel keinen merklichen Einfluss; in dem westlichsten Theile von Westphalen habe ich zwar Rücksicht darauf genommen, aber bloss zur Ehre der Rechnung, denn wirklich bringt es auch da nichts.
\begin{center}\textsc{Gauss} an \textsc{Schumacher}.\quad Göttingen, 9. December 1838.\end{center}
Die verlangte Formel ist folgende\textsuperscript{[*)]}:
\[
\begin{aligned}
\varphi={}&\psi+(6ff+48f^4+426f^6+4080f^8\ldots)\sin2\psi\\
&\quad+(21f^4+336f^6+4264f^8\ldots)\sin4\psi\\
&\qquad+\left(\frac{302}3f^6+2416f^8\ldots\right)\sin6\psi\\
&\qquad\quad+\left(\frac{1097}2f^8\ldots\right)\sin8\psi\\
&\qquad\qquad\text{etc.}
\end{aligned}
\]
\par\noindent\rule{4cm}{.4pt}\par
{\footnotesize\noindent\hspace*{\parindent}[*) Vergl. Art. 25, S. 179.]\par}
\noindent{\small IX.}\hfill28\hspace*{.12\linewidth}\par

% Printed218; principalPDF225.
\par\clearpage
\noindent218\hfill{\footnotesize\textsc{Briefwechsel.\quad Über die Formeln für die Hannoversche Landesvermessung.}}\par\medskip
Es hat mich, da ich mich seit ziemlich vielen Jahren mit diesen Dingen nicht beschäftigt habe, erst viel Suchens gekostet, bis ich die Formel wieder aufgefunden habe, und nachher neuen Suchens oder vielmehr Rechnens, um die Bedeutung des \(f\) zu ermitteln, die dem Blatt nicht beigeschrieben war. Es ist aber \(f=\frac14e\), wo \(e\) die Excentricität der erzeugenden Ellipse bedeutet; oder, wenn die Abplattung \(\frac1n\) ist, d. i. \(\frac{\text{Erdaxe}}{\text{Durchm. d. Äquators}}=\frac{n-1}n\), so ist \(f=\frac{\surd(2n-1)}{4n}\). Ich habe dieses \(f\) deswegen gebraucht, weil, wenn ich \(e\) gebraucht hätte, sämmtliche Coefficienten Brüche geworden wären, da sie hier grösstentheils ganze Zahlen sind. Die numerische Rechnung hatte ich durchgehends mit Logarithmen auf 10 Decimalen geführt.
\begin{center}\rule{3cm}{.4pt}\end{center}
\bigskip
\begin{center}\textsc{Bemerkungen.}\end{center}
\begingroup\footnotesize
Die ersten beiden und der letzte der vorstehenden Briefe sind nach den im \textsc{Gauss}-Archiv befindlichen Originalen abgedruckt, während für die Briefe vom 17. Mai und 25. Juni 1831 nach den Originalen angefertigte Copien benutzt wurden. Im Abdruck sind einige Schreib- und kleinere Rechenfehler berichtigt worden. Der Brief vom 9. December 1838 ist die Antwort auf eine Bitte \textsc{Schumachers}, ihm die Ableitung der im Briefe vom 18. April 1830 gegebenen numerischen Formel für \(\varphi-\psi\) mitzutheilen. Die \textsc{Gauss}sche Angabe
\[
\varphi=\psi+520''{,}469\,3364\sin2\psi\ \text{etc.}
\]
\noindent enthält jedoch einen Rechenfehler, wie auch das noch vorhandene Blatt zeigt, das \textsc{Gauss} zur Rechnung benutzt hatte; es muss heissen, wie S. 206 angegeben ist:
\[
\varphi=\psi+520''{,}463\,3364\sin2\psi\ \text{etc.}
\]
\noindent Dadurch wird auch die letzte Stelle in der dort gegebenen Tabelle für \(\varphi-\psi\) unsicher; ferner erhält man für Göttingen \(\psi^\circ=51^\circ23'20''{,}6082\), S. 207, an Stelle von \(\ldots20''{,}6024\), wie es im Original heisst.
\par\noindent\hfill\textsc{Krüger, Börsch.}\par
\endgroup
\vspace{1.8cm}\begin{center}\rule{4cm}{.4pt}\end{center}

% Unnumbered title219; principalPDF226.
\par\clearpage\thispagestyle{empty}
\null\vspace*{.23\textheight}
\begin{center}{\Large TRIGONOMETRISCHE PUNKTBESTIMMUNG.}\end{center}
\vfill\noindent\hfill28\textsuperscript{*}\hspace*{.15\linewidth}\par\vspace*{.14\textheight}

% Blank220; principalPDF227.
\par\clearpage\null\thispagestyle{empty}

% Unnumbered opening221; principalPDF228.
\par\clearpage\thispagestyle{empty}
\vspace*{2.8\baselineskip}
\begin{center}{\Large\textls[160]{NACHLASS.}}\end{center}
\vspace{2\baselineskip}
\begin{center}[1.]\\[5pt]Endresultat für den Ort eines Punktes in einer Ebene, der von drei bekannten\\aus angeschnitten ist.\end{center}
Es bedeuten \(10,20,30\) die drei beobachteten Richtungen [nach \(P\)] und \(\alpha,\beta,\gamma\) die entsprechenden Entfernungen.

Die drei einzelnen Resultate aus den Combinationen \(2-3,\ 1-3,\ 1-2\) seien \(A,B,C\), zugleich die Winkel des durch jene gebildeten Dreiecks; die ihnen gegenüber stehenden Seiten \(a,b,c\).
\begin{center}
\begin{tikzpicture}[x=1cm,y=1cm,line width=.6pt,font=\normalsize]
\coordinate (one) at (0,1.4);\coordinate (two) at (9.4,6.045);\coordinate (three) at (10.5,-.06);
\coordinate (B) at (4,2.8);\coordinate (C) at (7,3.85);\coordinate (A) at (5.25,2.25);\coordinate (P) at (5.44,2.86);
\draw (one)--(B)--(C);\draw (A)--(C)--(two);\draw (B)--(A)--(three);
\draw[line width=.4pt] (5.30155,3.25555)--(P)--(5.6573,2.62236) (P)--(5.1843,2.2789);
\foreach \p in {one,two,three,P}{\fill[white] (\p) circle(.045);\draw (\p) circle(.045);}
\node[above] at (one) {$1$};\node[right] at (two) {$2$};\node[right] at (three) {$3$};
\node[below] at (A) {$A$};\node[above left] at (B) {$B$};\node[right] at (C) {$C$};\node[above right] at (P) {$P$};
\node at (5.48,3.54) {$a$};\node at (6.46,2.92) {$b$};\node at (4.62,2.31) {$c$};
\node at (5.18,3.05) {$x$};\node at (5.53,2.55) {$y$};\node at (5.10,2.65) {$z$};
\end{tikzpicture}
\end{center}
